% TOP_PROBLEM: {"id": 3094, "title": "Edge-Unfolding Convex Polyhedra", "category_id": 6, "category": {"id": 6, "name": "geometry", "display_name": "Geometry", "description": "Euclidean and non-Euclidean geometry, geometric structures.", "slug": "geometry", "order_index": 6, "created_at": "2026-07-31T15:26:25.671Z"}, "status": "open", "problem_number": "OPG-60037", "set_id": 12}
% FIRST_POSTED: 2026-10-01
% ENTRY_KIND: statement_only
% UnsolvedMath ID 3094; source catalogue code OPG-60037
% !TeX program = lualatex
% Definition and sources only; this file is not an LLM solution attempt.
\documentclass[11pt,a4paper]{article}
\usepackage[margin=25mm]{geometry}
\usepackage{fontspec}
\setmainfont{lmroman10-regular.otf}[BoldFont=lmroman10-bold.otf,
 ItalicFont=lmroman10-italic.otf,BoldItalicFont=lmroman10-bolditalic.otf]
\usepackage{amsmath,amssymb,mathtools}
\usepackage{unicode-math}
\setmathfont{latinmodern-math.otf}
\AtBeginDocument{\let\setminus\smallsetminus}
\usepackage{xurl,hyperref}
\hypersetup{colorlinks=true,urlcolor=blue}
\setcounter{secnumdepth}{0}
\setlength{\parindent}{0pt}
\setlength{\parskip}{5pt}
\setlength{\emergencystretch}{3em}
\begin{document}
\section*{Edge-Unfolding Convex Polyhedra}\label{3094}
{\small UnsolvedMath ID 3094 · OPG-60037 · Geometry · OpenGarden}
\subsection{Definitions and mathematical statement}
Let $P\subset\mathbb R^3$ be a bounded convex polyhedron with nonempty
interior. Its boundary consists of finitely many planar polygonal faces,
meeting along its original edges and vertices. An edge-cut tree is a
spanning tree of the vertex-edge graph of $P$: it contains every vertex,
is connected, and has no cycle.

Cut the boundary surface along all edges of such a tree, duplicating the
cut edges on the resulting boundary. The remaining surface is a topological
disk made of the original polygonal faces, still attached along the uncut
edges. An edge unfolding places this disk into $\mathbb R^2$, isometrically
on each face, with the two incident faces of each uncut edge attached
along their common image. It can be obtained by successively rotating
faces about uncut edges until the surface is flat.

The conjecture states that every $P$ admits a choice of edge-cut tree
whose unfolded faces have pairwise disjoint interiors. This is the usual
nonoverlap condition for a net; contacts along face boundaries are allowed.
The flat result must be one connected piece containing every face exactly
once. Cutting through the interior of a face, adding new vertices to
supply extra cuts, or separating the faces into several independent pieces
is not allowed.

The problem is about existence of the final planar net. It does not
require that a physical motion from the original polyhedron to that net
avoid every intermediate collision. Convexity is imposed on the original
three-dimensional solid, not on the planar net, which may be nonconvex.
\subsection{Short English statement}
Can every convex polyhedron be cut only along an edge spanning tree and flattened into a single net with no overlapping face interiors?
\subsection{Sources}
\begin{itemize}
\item[S1] ULAM.AI and the UnsolvedMath Contributors, supplied problems.json, record 3094 (OPG-60037), OpenGarden collection. Catalogue identity and source transcription; CC BY 4.0.
\url{https://huggingface.co/datasets/ulamai/UnsolvedMath}.
\item[S2] Open Problem Garden, Edge-Unfolding Convex Polyhedra. Original problem and discussion; the mathematical formulation below is expanded from the supplied source record.
\url{http://www.openproblemgarden.org/op/edge_unfolding_convex_polyhedra}.
\end{itemize}
\end{document}
